\section{Introduzione al Problema}\label{introduzione-al-problema}

Nel lavoro di tesi, vogliamo gestire una competizione CTF.

Di solito, in una competizione CTF, uno o piu\textquotesingle{} utenti
(eventualmente organizzati in team) si sfidano per risolvere
un\textquotesingle insieme di challenge. Generalmente una challenge vuol
dire: recuperare una \textbf{flag}.

Ci sono vari tipi di competizioni CTF, tra cui:

\begin{itemize}
\tightlist
\item
  \textbf{Jeopardy}: gli utenti devono risolvere vari problemi, la cui
  flag e\textquotesingle{} di solito una stringa del tipo
  \passthrough{\lstinline!\{hai\_trovato\_la\_flag\}!}
\item
  \textbf{Attacco e Difesa}: e\textquotesingle{} un tipo di competizione
  dove ogni team, deve recuperare una o piu flag attaccando i sistemi
  dei team avversare, mentre ognuno si occupa di difendere il proprio
  sistema.
\end{itemize}

In particolare, con il nostro lavoro, vogliamo gestire uno o
piu\textquotesingle{} sistemi vulnerabili ad un qualche tipo di attacco.

La flag si ottiene difendendo il sistema, ossia \emph{"patchando"} la
vulnerabilita. Dunque la flag non e\textquotesingle{} piu una stringa
come nelle \textbf{Jeopardy}, ma piuttosto e\textquotesingle: uno stato
che deve raggiungere il sistema target mediante il susseguirsi di azioni
che il partecipante applica quest\textquotesingle ultimo.

Nel lavoro si tiene conte del fatto che il sistema vulnerabile, spesso,
potrebbe essere costituito da uno o piu host/nodi che comunicano tra di
loro. Per esempio si condideri il seguente scenario: un server
vulnerabile ha bisogno di comunicare con database in esecuzione su un
nodo differente da quello del server.
